\chapter{效率修正、信号提取与系统误差}

\section{接受度与效率分解}

参考公开 CMS quarkonium 分析的一般写法，本章建议把修正因子拆分为接受度、重建效率、鉴别效率、
触发效率以及必要的事件级修正因子。
若本分析最终不报告物理截面，也应至少明确这些修正因子在候选产额解释中的作用。

\section{信号提取策略}

本节用于说明最终结果是如何从数据中提取的。
根据后续分析实现，可在这里填入：

\begin{itemize}
  \item 质量分布或多维分布的拟合变量；
  \item 信号模型与组合背景模型；
  \item 参数固定项与自由项的来源；
  \item 交叉检验方法，例如旁带检验、替代模型或 toy study。
\end{itemize}

\section{系统误差框架}

系统误差部分应围绕分析链条逐项组织，而不是简单堆叠来源名称。
推荐的归类方式包括：

\begin{itemize}
  \item 触发、重建和鉴别效率；
  \item 质量模型和背景模型；
  \item 候选选择和优选规则；
  \item MC 建模与重加权；
  \item 积分亮度或归一化相关误差。
\end{itemize}

在最终数值尚未确定之前，这一章只保留方法框架和待填位置，不写假定百分比。
